\section*{Bab \ref{ch:g9:weight-vs-mass} --- Berat lawan Massa: Mengukur Gaya}

\begin{solution}{exo:g9:weight-vs-mass:1}
Newton (\unit{N}), yang diukur dengan \omterm{def:g9:weight-vs-mass:dynamometer}{dinamometer}. Jangkarnya: satu
\omterm{def:g9:weight-vs-mass:newton}{newton} --- \omterm{def:g4:weight-and-mass:weight}{berat} sebutir apel kecil; seratus --- sebuah jabat tangan
yang mantap (atau \omterm{def:g4:weight-and-mass:weight}{berat} sebuah koper \qty{10}{kg}).
\end{solution}

\begin{solution}{exo:g9:weight-vs-mass:2}
$P = m \times g$: \omterm{def:g4:weight-and-mass:weight}{berat} dalam \omterm{def:g9:weight-vs-mass:newton}{newton}, \omterm{def:g3:measuring-mass:mass}{massa} dalam \omterm{def:g3:measuring-mass:units}{kilogram}, dan
$g$ --- \omterm{prop:g9:weight-vs-mass:formula}{kuat gravitasi} tempat itu --- yaitu berapa \omterm{def:g9:weight-vs-mass:newton}{newton} tarikan yang
diterima tiap \omterm{def:g3:measuring-mass:units}{kilogramnya} di sana (\qty{9.8}{N/kg} di Bumi).
\end{solution}

\begin{solution}{exo:g9:weight-vs-mass:3}
Anjingnya: $25 \times 9.8 = \qty{245}{N}$. Rotinya: $0.7 \times 9.8
\approx \qty{6.9}{N}$. Seorang pembaca \qty{55}{kg}: kira-kira
\qty{539}{N}.
\end{solution}

\begin{solution}{exo:g9:weight-vs-mass:4}
$m = 14.7 \div 9.8 = \qty{1.5}{kg}$.
\end{solution}

\begin{solution}{exo:g9:weight-vs-mass:5}
Kemiringannya adalah $g$ setempat --- \omterm{def:g9:weight-vs-mass:newton}{newton} tiap \omterm{def:g3:measuring-mass:units}{kilogram}. Di \omterm{def:g2:moon-in-the-sky:moon}{Bulan}
garisnya tetap lurus melewati titik asalnya tetapi melorot menjadi
kemiringan $1.6$ yang landai: kesebandingan yang sama, dunia yang
lebih lemah.
\end{solution}

\begin{solution}{exo:g9:weight-vs-mass:6}
Lajur \omterm{def:g3:measuring-mass:mass}{massanya} tidak pernah berubah --- \qty{10}{kg} zat bepergian
utuh. Lajur $g$ dan lajur \omterm{def:g4:weight-and-mass:weight}{beratnya} berubah menurut dunianya: \omterm{def:g4:weight-and-mass:weight}{berat}
adalah pungutan setempat, $m \times g$.
\end{solution}

\begin{solution}{exo:g9:weight-vs-mass:7}
\omterm{def:g3:levers-and-scales:scale}{Neracanya} membandingkan $m_1 g$ dengan $m_2 g$: $g$ mengalikan kedua
piringnya lalu meniadakan diri --- \omterm{def:g3:measuring-mass:mass}{massa}, dengan benar, di mana pun.
Timbangan kamar mandinya \omterm{def:g6:measurement-in-science:measuring}{mengukur} $P = m g_{\text{di sini}}$ tetapi
membagi dengan $g_{\text{Bumi}}$ untuk mencetak \omterm{def:g3:measuring-mass:units}{kilogram}: di \omterm{def:g2:moon-in-the-sky:moon}{Bulan} ia
mencetak $m \times 1.6/9.8$ --- seperenam kebenarannya.
\end{solution}

\begin{solution}{exo:g9:weight-vs-mass:8}
\omterm{def:g4:weight-and-mass:weight}{Beratnya}: \qty{19.6}{N} ke bawah --- sebuah anak panah sepanjang
kira-kira \qty{3.9}{cm} pada lima \omterm{def:g9:weight-vs-mass:newton}{newton} tiap \omterm{def:g2:measuring-length:units}{sentimeter}, dari bukunya
ke bawah. Penyangganya: anak panah \qty{3.9}{cm} yang sama ke atas.
Panjang yang sama, arah yang berlawanan: anggarannya seimbang,
bukunya diam.
\end{solution}

\begin{solution}{exo:g9:weight-vs-mass:9}
Pesawatnya harus membawa, mengangkat, dan mengerem \emph{zat} tasnya
--- yaitu \omterm{def:g3:measuring-mass:mass}{massanya} --- dan \qty{23}{kg} zat tetap \qty{23}{kg} di dunia
mana pun. Meja bulannya memerlukan sebuah \omterm{def:g3:levers-and-scales:scale}{neraca} dengan \omterm{def:g3:measuring-mass:mass}{massa}
bertanda: satu-satunya timbangan yang selamat melewati perubahan $g$.
\end{solution}

\begin{solution}{exo:g9:weight-vs-mass:10}
$m = 111 \div 3.7 = \qty{30}{kg}$; di Bumi ia berbobot $30
\times 9.8 = \qty{294}{N}$.
\end{solution}

\begin{solution}{exo:g9:weight-vs-mass:11}
Ia membaca terlalu tinggi: ketika menggantung, pegasnya juga membawa
\omterm{def:g4:weight-and-mass:weight}{berat} dirinya sendiri (dan \omterm{def:g4:weight-and-mass:weight}{berat} kaitnya), yang tidak pernah dimuat
nol ratanya --- alatnya menambahkan kelebihan yang tetap. Nolkan ia
ketika menggantung, dan kelebihannya pun dikurangkan sebelum bebannya
tiba.
\end{solution}

\begin{solution}{exo:g9:weight-vs-mass:12}
Perlengkapannya: \omterm{def:g9:weight-vs-mass:dynamometer}{dinamometer}, seperangkat \omterm{def:g3:measuring-mass:mass}{massa} bertanda ($1$ sampai
\qty{5}{kg}), sebuah kait. Gantungkan tiap \omterm{def:g3:measuring-mass:mass}{massanya}, catatlah pasangan
($m$, $P$)-nya, lalu gambarkan $P$ terhadap $m$: titiknya sebaris
melewati titik asalnya, dan kemiringannya --- naiknya dalam \omterm{def:g9:weight-vs-mass:newton}{newton}
dibagi mendatarnya dalam \omterm{def:g3:measuring-mass:units}{kilogram} --- adalah $g \approx \qty{9.8}{N/kg}$.
Di \omterm{def:g2:moon-in-the-sky:moon}{Bulan}: tata cara yang serupa, kelurusan yang serupa, kemiringan
baru $1.6$ --- caranya \omterm{def:g6:measurement-in-science:measuring}{mengukur} dunia mana pun yang dipijaknya; hanya
bilangannya yang setempat.
\end{solution}

\begin{solution}{pb:g9:weight-vs-mass:1}
\textbf{1.} Bumi $70 \times 9.8 = \qty{686}{N}$; \omterm{def:g2:moon-in-the-sky:moon}{Bulan}
\qty{112}{N}; Mars \qty{259}{N}; geladak Jupiter
\qty{1736}{N}.
\textbf{2.} Benar di Bumi: \qty{70}{kg}. \omterm{def:g2:moon-in-the-sky:moon}{Bulan}: $112 \div
9.8 \approx \qty{11.4}{kg}$; Mars: $259 \div 9.8 \approx
\qty{26.4}{kg}$; Jupiter: $1736 \div 9.8 \approx
\qty{177}{kg}$ --- sanjungan dan fitnah menurut rumus.
\textbf{3.} \qty{70}{kg} di tiap pelabuhannya: perbandingan \omterm{def:g3:levers-and-scales:scale}{neracanya}
meniadakan $g$ setempatnya.
\textbf{4.} Hanya geladak Jupiter: \qty{1736}{N} mendekati batasnya
tetapi tetap di bawahnya --- tidak ada pelabuhan yang melanggar
\qty{2000}{N} bagi pelancong ini (seorang rekan yang lebih berat
seberat \qty{90}{kg} akan melanggarnya di sana:
$90 \times 24.8 = \qty{2232}{N}$).
\textbf{5.} Keberangkatannya: $40 \times 9.8 = \qty{392}{N}$;
kedatangannya: $40 \times 3.7 = \qty{148}{N}$.
\textbf{6.} $m = 200/g$: Bumi \qty{20.4}{kg}; \omterm{def:g2:moon-in-the-sky:moon}{Bulan}
\qty{125}{kg}; Mars \qty{54}{kg}; Jupiter \qty{8.1}{kg} ---
\omterm{def:g9:weight-vs-mass:newton}{newton} tetap milik serikatnya membeli \omterm{def:g3:measuring-mass:mass}{massa} yang sangat berbeda-beda.
\textbf{7.} $m = 150 \div 1.6 = \qty{93.75}{kg}$ --- kira-kira
\qty{94}{kg} peti untuk bea cukainya.
\textbf{8.} $m = 294 \div 9.8 = \qty{30}{kg}$; \omterm{def:g4:weight-and-mass:weight}{beratnya} di \omterm{def:g2:moon-in-the-sky:moon}{Bulan}
$30 \times 1.6 = \qty{48}{N}$.
\textbf{9.} Otot dan \omterm{def:g3:measuring-mass:mass}{massanya} pulang tanpa berubah; hanya lawannya
yang berubah. Di \omterm{def:g2:moon-in-the-sky:moon}{Bulan} tiap \omterm{def:g3:measuring-mass:units}{kilogramnya} memakan \qty{1.6}{N} untuk
diangkat; di rumahnya menagih \qty{9.8}{N} lagi. Tambahannya adalah
potongan harga \omterm{def:g2:moon-in-the-sky:moon}{Bulan}, bukan pertumbuhan tubuhnya --- $P = m g$, dengan
$m$ yang setia dan $g$ yang berubah-ubah.
\textbf{10.} $0.2 \times 24.8 \approx \qty{5}{N}$ --- ia bukan
landasan tempa: di Bumi pengocoknya berbobot \qty{2}{N}, sehingga
mengocok di geladaknya dua setengah kali lebih berat, melelahkan bagi
pergelangan tangan tetapi hampir bukan pekerjaan pengecoran logam.
Keluhannya: dikabulkan sebagian.
\textbf{11.} Senewton cokelat adalah $m = 1/g$: Bumi kira-kira
\qty{102}{g}; \omterm{def:g2:moon-in-the-sky:moon}{Bulan} \qty{625}{g}; Jupiter \qty{40}{g}. Belilah
newtonmu di \omterm{def:g2:moon-in-the-sky:moon}{Bulan}.
\textbf{12.} Misalnya: ``\omterm{def:g3:measuring-mass:mass}{Massa} bepergian bersamamu; \omterm{def:g4:weight-and-mass:weight}{berat} adalah
pungutan tiap dunia atasnya --- kemasilah \omterm{def:g3:measuring-mass:units}{kilogram}, lalu biarkan
pelabuhannya menagih \omterm{def:g9:weight-vs-mass:newton}{newtonnya}.''
\end{solution}
